\documentclass[11pt]{article}

\usepackage[margin=1in]{geometry}
\usepackage{amsmath,amssymb,amsthm,mathtools}
\usepackage{booktabs,enumitem}
\usepackage{algorithm,algpseudocode}
\usepackage{microtype}
\usepackage{needspace}
\usepackage[hidelinks]{hyperref}

\newcommand{\studentname}{Xilang Chen}
\newcommand{\studentid}{524031910668}
\DeclarePairedDelimiter{\floor}{\lfloor}{\rfloor}
\DeclarePairedDelimiter{\ceil}{\lceil}{\rceil}
\newcommand{\dist}{\operatorname{dist}}
\newcommand{\key}{\operatorname{key}}
\newcommand{\pos}{\operatorname{pos}}
\newcommand{\MST}{\operatorname{MST}}
\newcommand{\exercise}[1]{\Needspace{6\baselineskip}\section*{Exercise #1}}
\setlength{\parindent}{0pt}
\setlength{\parskip}{0.45em}
\setlist{itemsep=0.25em,topsep=0.3em}
\raggedbottom

\title{SE3308 Algorithm Design I --- Assignment 3\\[0.3em]
\large Exercises 4.11, 4.12, 4.16, 5.4, 5.6, and 5.10}
\author{\studentname\quad\studentid}
\date{}

\begin{document}
\maketitle

Write $n=|V|$ and $m=|E|$. Graphs use adjacency lists and, as usual, have no
parallel edges, so $m=O(n^2)$. Undirected graphs have no self-loops.
Edge-weight arithmetic and comparisons take constant time. An unreachable
distance is $+\infty$; $\log$ means $\log_2$.

\exercise{4.11}

\textbf{Algorithm.} For each vertex $s$, run Dijkstra's algorithm from $s$,
using an unsorted array for the priority queue. For every incoming edge
$(u,s)$, consider the shortest path from $s$ to $u$ followed by $(u,s)$.
Return
\[
  \boxed{L=\min_{(u,v)\in E}\bigl(\dist_G(v,u)+\ell(u,v)\bigr).}
\]
If every candidate is infinite, report that $G$ is acyclic.

\begin{algorithm}[H]
\caption{Shortest directed cycle}
\begin{algorithmic}[1]
\Function{ShortestCycle}{$G,\ell$}
  \State $L\gets+\infty$
  \For{each $s\in V$}
    \State $d\gets\Call{ArrayDijkstra}{G,\ell,s}$
    \For{each incoming edge $(u,s)$}
      \State $L\gets\min\{L,d[u]+\ell(u,s)\}$
    \EndFor
  \EndFor
  \If{$L=+\infty$} \State \Return \textnormal{acyclic} \EndIf
  \State \Return $L$
\EndFunction
\end{algorithmic}
\end{algorithm}

\textbf{Correctness.} Positive edge lengths imply that a shortest path between
distinct vertices can be taken to be simple: removing a repeated-vertex
segment strictly shortens a walk. Thus, whenever $\dist_G(v,u)$ is finite,
a shortest $v$--$u$ path followed by $(u,v)$ is a directed simple cycle.
The path cannot already use $(u,v)$, since that would revisit its starting
vertex $v$. If self-loops are permitted, $(v,v)$ itself is a cycle and the
same formula uses $\dist_G(v,v)=0$.

\Needspace{7\baselineskip}
Let $\mathrm{OPT}$ be the length of a shortest cycle, when one exists.
Every finite candidate is the length of a cycle, so $L\ge\mathrm{OPT}$.
Conversely, choose an edge $(u,v)$ of a shortest cycle $C$. Removing that
edge leaves a $v$--$u$ path, and hence
\[
  L\le \dist_G(v,u)+\ell(u,v)\le \ell(C)=\mathrm{OPT}.
\]
Therefore $L=\mathrm{OPT}$. A finite candidate exists if and only if a cycle
exists, which also proves the acyclicity report correct.

\textbf{Running time.} Array-based Dijkstra performs at most $n$ linear
minimum searches and scans every edge once: $O(n^2+m)=O(n^2)$ per source.
Build incoming adjacency lists once in $O(n+m)$ time; their scans across
all sources total $O(m)$. The total is
\[
  O\bigl(n(n^2+m)+n+m\bigr)=\boxed{O(n^3)}.
\]
Only one source's distance array is needed at a time.

\exercise{4.12}

Let the specified edge be $e=\{a,b\}$, and let $G'=G-e$.

\textbf{Algorithm.} Run array-based Dijkstra from $a$ in $G'$, ignoring
$e$ in both endpoints' adjacency lists. If $b$ is unreachable, report
that no cycle contains $e$ (equivalently, return $+\infty$). Otherwise return
\[
  \boxed{\ell(e)+\dist_{G-e}(a,b).}
\]

\begin{algorithm}[H]
\caption{Shortest cycle through a specified undirected edge}
\begin{algorithmic}[1]
\Function{ShortestCycleThroughEdge}{$G,\ell,e=\{a,b\}$}
  \State $d\gets\Call{ArrayDijkstra}{G-e,\ell,a}$
  \If{$d[b]=+\infty$}
    \State \Return \textnormal{no cycle contains $e$}
  \EndIf
  \State \Return $\ell(e)+d[b]$
\EndFunction
\end{algorithmic}
\end{algorithm}

\textbf{Correctness.} Removing $e$ from any simple cycle containing $e$
leaves an $a$--$b$ path in $G-e$. Consequently every such cycle has length
at least $\ell(e)+\dist_{G-e}(a,b)$.

Conversely, if $b$ is reachable, take a shortest $a$--$b$ path $P$ in
$G-e$. By positivity, $P$ is simple. It avoids $e$, so $P\cup\{e\}$ is a
simple cycle of exactly the returned length. In a simple undirected
graph, $P$ has at least two edges. If $b$ is unreachable, the first
direction shows that no cycle through $e$ can exist.

\textbf{Running time.} One array-based Dijkstra run costs
$O(n^2+m)=\boxed{O(n^2)}$.
Skipping $e$ requires only constant extra work per adjacency entry;
there is no need to copy the graph.

\textbf{Why remove $e$?} In the original graph, a shortest $a$--$b$ path
could be $e$ itself. Following it and then traversing $e$ backwards is
not a simple cycle.

\exercise{4.16}

All arrays are indexed from $1$, and complete trees are stored in
level order, left to right. We call the final subpart \textbf{(d)};
it is printed as a second \textbf{(a)} in the exercise.

\textbf{(a) Binary-tree indices.}
The root is at index $1$. The two child slots of each successive parent
occupy the pairs
\[
 (2,3),\ (4,5),\ (6,7),\ \ldots.
\]
Thus the children of index $j$ are $2j$ and $2j+1$, retaining only
indices at most $n$. Conversely, every non-root index $j\ge2$ is
either $2p$ or $2p+1$ for its parent $p$, so
\[
  \boxed{\operatorname{parent}(j)=\floor{j/2},\qquad j\ge2.}
\]
The root has no parent.

\textbf{(b) $d$-ary-tree indices.}
Assume $d\ge2$. The root's child slots start at index $2$.
The preceding $j-1$ parents account for $d(j-1)$ child slots, so the
$r$-th child of parent $j$ has index
\[
  \boxed{\operatorname{child}_r(j)=d(j-1)+r+1,\qquad 1\le r\le d.}
\]
Discard indices exceeding $n$.
A non-root index $q$ belongs to the block of its parent $p$ exactly when
\[
 d(p-1)+2\le q\le d(p-1)+d+1.
\]
Therefore
\[
  \boxed{\operatorname{parent}(q)=
    \floor{\frac{q-2}{d}}+1
    =\ceil{\frac{q-1}{d}},\qquad q\ge2.}
\]
Putting $d=2$ recovers part (a).

\Needspace{10\baselineskip}
\textbf{(c) Linear-time makeheap and a worst-case input.}

\emph{Why bottom-up heap construction works.}
When \textsc{siftdown} is called at position $i$, its child subtrees are
already min-heaps: every child has index greater than $i$ and was
processed earlier. Moving the smaller child up until the saved item
fits restores the heap property throughout the subtree rooted at $i$.
Induction over decreasing $i$ proves that the final array is a min-heap.

\emph{Running time.}
The height $H_i$ of the subtree rooted at $i$ is
\[
 H_i=\floor{\log(n/i)}.
\]
Indeed, the leftmost descendant $r$ levels below $i$ has index $2^r i$,
so such a level exists exactly when $2^r i\le n$.
The call at $i$ costs $O(1+H_i)$. Using the individual heights avoids
the loose upper bound $O(n\log n)$. Counting the total levels gives
\[
 \begin{aligned}
 \sum_{i=1}^n H_i
 &=\sum_{r=1}^{\floor{\log n}}
       \bigl|\{i:2^r i\le n\}\bigr|\\
 &=\sum_{r=1}^{\floor{\log n}}\floor{\frac{n}{2^r}}
 \le n\sum_{r=1}^{\infty}2^{-r}=n.
 \end{aligned}
\]
Thus filling the array and all calls together take
$O(n+\sum_{i=1}^n H_i)=\boxed{O(n)}$.
This also justifies the hint: since
$\log(n/i)<H_i+1$, we have $\sum_{i=1}^n\log(n/i)<2n$.
The additive $O(n)$ accounts for initialization and constant work at
leaves. Initialization alone gives an $\Omega(n)$ lower bound.

\Needspace{10\baselineskip}
\emph{Worst-case input family.}
For a perfect binary tree with $n=2^{h+1}-1$ nodes, put distinct keys in
strictly decreasing array order:
\[
  h[1..n]=[n,n-1,\ldots,1].
\]
At the call for index $i$, the saved key is larger than every key in
its descendant subtrees. Earlier calls have only rearranged keys within
those subtrees. It therefore moves down to a leaf; because the tree is
perfect, every leaf in that subtree is $H_i$ levels below $i$.
Every call attains its maximum possible downward distance.
The number of downward moves is
\[
  \sum_i H_i
  =\sum_{r=1}^{h}(2^{h+1-r}-1)
  =2^{h+1}-h-2=n-h-1.
\]
This is a worst-case family with $\Theta(n)$ work, matching the upper
bound. For an incomplete last level, decreasing order need not
maximize each individual descent, but the worst-case asymptotic
bound remains $\Theta(n)$.

\Needspace{12\baselineskip}
\textbf{(d) Adapting the pseudocode to a $d$-ary heap.}

Use the parent and child formulas from (b). In particular, replace every
parent computation in \textsc{bubbleup} by
\[
   p\gets\floor{\frac{i-2}{d}}+1\qquad(i>1).
\]
The root case must be checked before accessing its parent.
Replace \textsc{minchild} by the following procedure.

\begin{algorithm}[H]
\caption{Smallest child in a $d$-ary heap}
\begin{algorithmic}[1]
\Function{MinChild}{$h,i,d$}
  \State $a\gets d(i-1)+2$
  \If{$a>|h|$} \State \Return $0$ \EndIf
  \State $b\gets\min\{|h|,a+d-1\}$
  \State \Return an index in $[a,b]$ minimizing $\key(h[j])$
\EndFunction
\end{algorithmic}
\end{algorithm}

The control flow of \textsc{siftdown} is unchanged; it now selects among
up to $d$ children. The same \textsc{insert}, \textsc{decreasekey}, and
bottom-up \textsc{makeheap} schemes apply. The latter may skip leaves
and start at $\floor{(n-2)/d}+1$ when $n\ge2$.

Maintain the inverse map $h^{-1}(x)=\pos[x]$: whenever an item $x$ is
placed in slot $j$, also set $\pos[x]\gets j$. This preserves the
constant-time lookup required by \textsc{decreasekey}.

Two details also apply to the binary version:
\begin{itemize}
\item Figure 4.16 prints $\ceil{i/2}$ in \textsc{bubbleup}; for a
  $1$-indexed binary heap the correct parent is $\floor{i/2}$.
\item \textsc{deletemin} must actually remove the last slot. Save the
  root and last item, reduce the logical heap size by one, and sift the
  saved last item from the root only if the new heap is nonempty.
  Remove the deleted item's inverse-map entry. Array capacity is
  distinct from logical heap size.
\end{itemize}

\exercise{5.4}

\textbf{Claim.} An undirected graph with $n$ vertices and $k$ connected
components has at least $n-k$ edges.

\begin{proof}
Start with the same $n$ vertices and no edges, hence $n$ components.
Insert the graph's $m$ edges one at a time. Each inserted edge either
joins two components and reduces their number by one, or has both
endpoints in the same component and leaves their number unchanged.
After $m$ insertions there are therefore at least $n-m$ components.
The final number is $k$, so
\[
  k\ge n-m,\qquad\text{equivalently}\qquad \boxed{m\ge n-k}.\qedhere
\]
\end{proof}

The bound is tight: a forest with $k$ components has exactly $n-k$
edges. Equivalently, if component $i$ has $n_i$ vertices, its spanning
tree contributes $n_i-1$ edges, and
$\sum_{i=1}^k(n_i-1)=n-k$.

\exercise{5.6}

We assume $G$ is connected, as is necessary for a spanning tree to exist.
For a disconnected graph, the corresponding conclusion is uniqueness
of the minimum spanning forest.

\textbf{Proof by exchange.}
Suppose two distinct MSTs $T$ and $T'$ exist. Let $e$ be the lightest
edge in their symmetric difference
\[
  E(T)\triangle E(T')
   =\bigl(E(T)\setminus E(T')\bigr)
     \cup\bigl(E(T')\setminus E(T)\bigr).
\]
Interchange the names of the trees if necessary so that
$e\in E(T)\setminus E(T')$.

Adding $e$ to $T'$ creates a unique cycle $C$.
Some edge $f$ of $C$ lies in $E(T')\setminus E(T)$: otherwise every edge
of $C$, including $e$, would belong to $T$, contradicting that $T$ is a tree.
Both $e$ and $f$ are in the symmetric difference. Since $e$ is its
lightest edge and all edge weights are distinct,
\[
  w(e)<w(f).
\]
Deleting $f$ from $T'+e$ breaks the unique cycle, leaving a spanning tree
$T''$ with
\[
  w(T'')=w(T')+w(e)-w(f)<w(T').
\]
This contradicts the minimality of $T'$. Hence
$\boxed{\text{the MST is unique.}}$ \hfill$\square$

\clearpage
\exercise{5.10}

Write $H=(V_H,E_H)$ with the weights inherited from $G$, and let
\[
  F=E(T)\cap E_H.
\]
We must prove that a \emph{single} MST of $H$ contains every edge of $F$.

Since $H$ is connected, it has an MST. Among all its MSTs choose one,
say $S$, maximizing $|E(S)\cap F|$. Suppose for contradiction that
some $e=\{u,v\}\in F$ does not belong to $S$.

Removing $e$ from $T$ partitions $V$ into two sets $A$ and $B$, with
$u\in A$ and $v\in B$. The $u$--$v$ path in $S$ must contain an edge
$f$ crossing this cut. As $e\notin E(S)$, we have $f\ne e$.
Moreover, $e$ is the \emph{only} edge of $T$ crossing the cut defined
by the two components of $T-e$. Thus $f\notin E(T)$, so in particular
$f\notin F$.

We claim $w(e)\le w(f)$. Otherwise $T-e+f$ would be a spanning tree of
$G$ of weight
\[
  w(T)-w(e)+w(f)<w(T),
\]
contradicting that $T$ is an MST.

Now insert $e$ into $S$. Its unique cycle consists of $e$ and the
$u$--$v$ path in $S$, so it contains $f$. Hence
\[
  S'=S+e-f
\]
is a spanning tree of $H$. Its weight satisfies
\[
  w(S')=w(S)+w(e)-w(f)\le w(S).
\]
Since $S$ is already minimum, equality holds and $S'$ is also an MST
of $H$. But $e\in F$ and $f\notin F$, so
\[
  |E(S')\cap F|=|E(S)\cap F|+1,
\]
contradicting the choice of $S$.

Therefore $F\subseteq E(S)$, proving
\[
  \boxed{T\cap H\text{ is contained in some MST of }H.}
\]
The proof allows equal edge weights and does not require $H$ to be
an induced subgraph. The intersection may be a disconnected forest;
the claim is that it can be extended to an MST of $H$.

\end{document}
